\documentclass[12pt]{article}
\usepackage[T1]{fontenc}
\usepackage[utf8]{inputenc}
\usepackage[polish]{babel}
\usepackage{lmodern}
\usepackage{geometry}
\geometry{a4paper,margin=2.5cm}
\usepackage{booktabs}
\usepackage{array}
\usepackage{url}
\usepackage{enumitem}
\usepackage{amsmath,amssymb}
\usepackage[hidelinks]{hyperref}
\setlength{\parindent}{1.25cm}
\setlength{\parskip}{4pt}
\renewcommand{\baselinestretch}{1.15}


\title{Raport projektowy\\Uczenie się ze wzmocnieniem}
\author{Vasyl Tesliak}
\date{2026/2027}

\begin{document}
\maketitle
\tableofcontents
\newpage

\section{Informacje podstawowe}

\begin{center}
\begin{tabular}{p{0.34\textwidth}p{0.56\textwidth}}
\toprule
Element & Informacja \\
\midrule
Imię i nazwisko & Vasyl Tesliak \\
Numer albumu & 181512 \\
Kierunek/specjalność & Informatyka/Sztuczna Inteligencja \\
Tytuł projektu & Agent uczący się optymalnej nawigacji w środowisku Grid World z ograniczonym zasobem energii \\
Link do repozytorium & \url{https://github.com/SkwiDwardClone/Reinforcement-learning-grid-world} \\
Data oddania & 18.11.2026 \\
\bottomrule
\end{tabular}
\end{center}

\section{Temat i cel projektu}

\subsection{Opis problemu}

Projekt dotyczy problemu nawigacji agenta w dyskretnym środowisku typu Grid World, przedstawionym jako dwuwymiarowa plansza podzielona na jednakowe pola. Agent rozpoczyna działanie w wyznaczonym polu startowym i ma za zadanie dotrzeć do pola docelowego. Przejście pomiędzy sąsiednimi polami wymaga wykonania jednej z dostępnych akcji ruchu.

Środowisko zawierać będzie statyczne przeszkody, których agent nie może przekraczać, pola niebezpieczne (pułapki) oraz pola umożliwiające odzyskanie części energii. Agent dysponuje ograniczonym zasobem energii, który zmniejsza się podczas poruszania. Wejście na pole regenerujące zwiększa poziom energii, z uwzględnieniem ustalonego limitu maksymalnego. Wyczerpanie energii uniemożliwia kontynuowanie epizodu.

Problem polega zatem nie tylko na znalezieniu krótkiej drogi do celu, lecz także na takim planowaniu kolejnych ruchów, aby agent nie utracił energii przed osiągnięciem celu i unikał niebezpiecznych pól. W pierwszej wersji projektu rozważane będzie środowisko o znanej, dyskretnej strukturze, a poziom trudności będzie zmieniany poprzez przygotowanie różnych map.

\subsection{Cel agenta}

Celem agenta jest dotarcie z pola startowego do pola docelowego przy możliwie małym koszcie poruszania się, unikaniu pułapek oraz efektywnym gospodarowaniu energią. Działanie agenta będzie uznawane za skuteczne, jeżeli w danym epizodzie osiągnie on pole docelowe przed wyczerpaniem energii i przed przekroczeniem maksymalnej liczby kroków.

Oprócz samego osiągnięcia celu istotna będzie jakość wyznaczonej strategii. Preferowane będzie rozwiązanie, które pozwala osiągać cel z wysoką częstotliwością, przy niewielkiej liczbie kroków i bez niepotrzebnego zużywania energii. Agent powinien nauczyć się uwzględniać zarówno bezpośrednie konsekwencje ruchu, jak i jego wpływ na możliwość wykonania kolejnych działań.

\subsection{Uzasadnienie zastosowania uczenia się ze wzmocnieniem}

Problem można sformułować jako zadanie uczenia się ze wzmocnieniem, ponieważ agent wielokrotnie podejmuje decyzje wchodzące w interakcję ze środowiskiem i otrzymuje informację zwrotną w postaci nagrody.

Decyzje są podejmowane sekwencyjnie: w każdym kroku agent wybiera kierunek ruchu na podstawie aktualnej pozycji i poziomu energii. Wykonany ruch zmienia jego położenie, może zmienić poziom energii oraz wpływa na zbiór możliwości dostępnych w kolejnych krokach. Przykładowo, wybór krótszej trasy może prowadzić przez pole niebezpieczne, natomiast dłuższa trasa może umożliwić uzupełnienie energii.

Istotny jest wynik długoterminowy, ponieważ pojedynczy ruch dający krótkoterminową korzyść nie musi prowadzić do osiągnięcia celu. Agent powinien nauczyć się strategii maksymalizującej sumę nagród w całym epizodzie, a nie wybierać wyłącznie ruchy o najlepszej natychmiastowej ocenie.

Zastosowanie agenta uczącego się jest uzasadnione możliwością porównania strategii wypracowanej na podstawie doświadczenia z prostym agentem bazowym, np. wybierającym losowe akcje. Uczenie ze wzmocnieniem pozwoli zbadać, czy agent jest w stanie samodzielnie wypracować skuteczną politykę nawigacji w różnych wariantach planszy.

\section{Formalizacja problemu RL/MDP}

\subsection{Agent i środowisko}

Agentem jest program podejmujący decyzje o kierunku ruchu na planszy. Środowiskiem jest dyskretna plansza Grid World zawierająca pole startowe, pole docelowe, przeszkody, pułapki oraz pola regenerujące energię. Środowisko przechowuje aktualną pozycję agenta, jego poziom energii oraz informację o rodzaju pól na planszy.

Jeden krok interakcji przebiega następująco:
\begin{enumerate}
    \item agent otrzymuje obserwację opisującą swoją pozycję i aktualny poziom energii;
    \item agent wybiera jedną z dostępnych akcji ruchu;
    \item środowisko sprawdza poprawność akcji i wyznacza nowe położenie agenta;
    \item środowisko aktualizuje poziom energii, uwzględniając koszt ruchu oraz ewentualne działanie pola regenerującego;
    \item środowisko oblicza nagrodę i sprawdza, czy wystąpił warunek zakończenia epizodu;
    \item agent otrzymuje nową obserwację, nagrodę oraz informację o zakończeniu epizodu.
\end{enumerate}

W podstawowym wariancie przejścia pomiędzy polami będą deterministyczne: poprawna akcja prowadzi do sąsiedniego pola w wybranym kierunku, o ile nie jest ono przeszkodą ani nie znajduje się poza planszą.

\subsection{Stan lub obserwacja}

W podstawowym wariancie agent będzie obserwował pełny stan potrzebny do podejmowania decyzji. Stan będzie składał się z aktualnej współrzędnej poziomej, aktualnej współrzędnej pionowej oraz bieżącego poziomu energii. Przy stałej mapie informacja o rozmieszczeniu przeszkód, pułapek i pól regenerujących jest częścią konfiguracji środowiska i będzie znana agentowi.

\begin{center}
\begin{tabular}{p{0.30\textwidth}p{0.60\textwidth}}
\toprule
Składowa stanu/obserwacji & Znaczenie \\
\midrule
Pozycja $x$ & Współrzędna pozioma pola zajmowanego przez agenta. \\
Pozycja $y$ & Współrzędna pionowa pola zajmowanego przez agenta. \\
Energia $e$ & Aktualny poziom energii agenta, ograniczony wartością maksymalną. \\
\bottomrule
\end{tabular}
\end{center}

Stan można zapisać jako $s_t=(x_t,y_t,e_t)$. Poziom energii będzie reprezentowany dyskretnie, co umożliwi zastosowanie tablicowej wersji algorytmu Q-learning. W przypadku zmiany mapy agent będzie trenowany i testowany z uwzględnieniem konfiguracji danej mapy.

\subsection{Akcje}

Zbiór akcji agenta obejmuje cztery kierunki ruchu. Wykonanie akcji powoduje próbę przejścia na sąsiednie pole. Jeżeli wskazane pole znajduje się poza planszą albo jest przeszkodą, ruch zostaje uznany za niedozwolony: agent pozostaje na dotychczasowym polu i otrzymuje karę. Takie rozwiązanie pozwala zachować stały zbiór czterech akcji w każdym stanie.

\begin{center}
\begin{tabular}{p{0.30\textwidth}p{0.60\textwidth}}
\toprule
Akcja & Opis \\
\midrule
GÓRA & Próba przejścia o jedno pole w górę. \\
DÓŁ & Próba przejścia o jedno pole w dół. \\
LEWO & Próba przejścia o jedno pole w lewo. \\
PRAWO & Próba przejścia o jedno pole w prawo. \\
\bottomrule
\end{tabular}
\end{center}

\subsection{Funkcja nagrody}

Wstępnie przyjmuje się następującą funkcję nagrody:
\[
r_t =
\begin{cases}
+100 & \text{osiągnięcie pola docelowego},\\
-30 & \text{wejście na pole pułapki},\\
-50 & \text{wyczerpanie energii},\\
-5 & \text{próba niedozwolonego ruchu},\\
-1 & \text{wykonanie zwykłego kroku}.
\end{cases}
\]

Nagroda za osiągnięcie celu jest wysoka, aby jednoznacznie premiować zakończenie zadania sukcesem. Kara za wejście na pułapkę oraz wyczerpanie energii ma zniechęcać do działań prowadzących do porażki. Niewielka kara za każdy zwykły krok zachęca do znajdowania krótszych tras i ogranicza niepotrzebne poruszanie się. Kara za niedozwolony ruch ma skłaniać agenta do uwzględniania granic planszy i położenia przeszkód.

Pola regenerujące energię zwiększają poziom energii o ustaloną wartość, nie przekraczając maksymalnego poziomu energii. W podstawowym wariancie samo uzupełnienie energii nie będzie powodowało dodatkowej nagrody, aby agent nie był zachęcany do wielokrotnego odwiedzania takich pól wyłącznie w celu zdobywania punktów. Wartości nagród są wstępne i mogą zostać skorygowane po przeprowadzeniu pierwszych testów środowiska.

\subsection{Warunki zakończenia epizodu}

Epizod zostaje zakończony w jednej z następujących sytuacji:
\begin{itemize}
    \item sukces: agent dociera do pola docelowego;
    \item porażka: agent wchodzi na pole pułapki;
    \item porażka: poziom energii agenta spada do zera;
    \item zakończenie z powodu limitu: agent osiąga maksymalną dopuszczalną liczbę kroków bez dotarcia do celu.
\end{itemize}

Osiągnięcie celu jest jedynym warunkiem oznaczającym sukces epizodu. Pozostałe wymienione sytuacje oznaczają jego niepowodzenie lub przerwanie.

\subsection{Założenia i ograniczenia środowiska}

W pierwszej wersji projektu przyjęte zostaną następujące założenia:
\begin{itemize}
    \item plansza ma postać skończonej, dwuwymiarowej siatki;
    \item położenie przeszkód, pułapek, pola startowego, celu i pól regenerujących jest stałe w trakcie pojedynczego epizodu;
    \item agent wykonuje jedną akcję ruchu w każdym kroku;
    \item przejścia pomiędzy polami są deterministyczne;
    \item energia jest zasobem dyskretnym, ma określony poziom maksymalny i zmienia się zgodnie z regułami środowiska;
    \item agent otrzymuje pełną informację o swojej pozycji i poziomie energii;
    \item w pierwszej wersji nie będą uwzględniane inni agenci ani nieprzewidywalne zdarzenia losowe.
\end{itemize}

Projekt będzie ograniczony do symulowanego środowiska, bez wykorzystania fizycznego robota. Planowane jest przygotowanie kilku map o różnym poziomie trudności. Wstępna formalizacja może zostać doprecyzowana po implementacji i testach, w szczególności w zakresie parametrów energii oraz wartości funkcji nagrody.
\end{document}